%
% motivation-definition.tex
%
% (c) 2026 Damien Flury, OST Ostschweizer Fachhochschule
%
% !TEX root = ../../buch.tex
% !TEX encoding = UTF-8
%
\section{Motivation und Definition
\label{hauptwert:section:motivation-definition}}
\kopfrechts{Teil 0}
\subsection{Uneigentliche Integrale}
\input{papers/hauptwert/fig/kehrwert.tex}

Man betrachte die Funktion
\begin{equation}
f\colon \mathbb{R}\setminus\{0\} \to \mathbb{R},
\qquad
f(x) = \frac{1}{x}
\end{equation}
und damit das Integral
\begin{equation}
\int_{-\infty}^{\infty} f(x)\,dx
= \int_{-\infty}^{\infty} \frac{1}{x}\,dx.
\end{equation}
Abbildung \ref{hauptwert:figure:kehrwert} zeigt den Graphen des Integranden.
Dieses Integral ist in zweifacher Hinsicht kritisch:
Der Integrand ist an der Stelle $x = 0$ nicht definiert
und hat dort eine Polstelle,
und die Integrationsgrenzen $\pm\infty$ sind selbst uneigentlich.
Gerade weil beide Arten von Divergenz gleichzeitig auftreten,
eignet sich dieses Integral als durchgehendes Beispiel.

\subsection{Ein erster Versuch}
Da der Integrand an der Stelle $x = 0$ nicht definiert und dort unbeschränkt
ist, handelt es sich nicht um ein Riemann-Integral.
Seine Bedeutung erhält der Ausdruck erst durch die Definition des
uneigentlichen Integrals, welche jede kritische Stelle durch einen eigenen,
unabhängigen Grenzwert ersetzt.
In einem ersten Versuch wird das Integral deshalb an der Polstelle in die
beiden Intervalle $(-\infty,0)$ und $(0,\infty)$ aufgeteilt:
\begin{equation}
\int_{-\infty}^{\infty} f(x)\,dx
= \lim_{\varepsilon_1\to 0^+} \int_{-\infty}^{-\varepsilon_1} f(x)\,dx
+ \lim_{\varepsilon_2\to 0^+} \int_{\varepsilon_2}^{\infty} f(x)\,dx.
\label{hauptwert:equation:aufteilung}
\end{equation}
Da auch die Integrationsgrenzen $\pm\infty$ uneigentlich sind,
muss jedes Teilintegral nochmals aufgeteilt
und an jedem uneigentlichen Ende durch einen eigenen Grenzwert ersetzt werden.
Das rechte Teilintegral wird dazu an einer beliebigen Zwischenstelle
$c > 0$ aufgeteilt und erhält am oberen Ende den Grenzwert $R \to \infty$,
das linke Teilintegral einen davon unabhängigen Grenzwert $R' \to -\infty$:
\begin{equation}
\lim_{\varepsilon_2\to 0^+} \int_{\varepsilon_2}^{c} \frac{1}{x}\,dx
+ \lim_{R\to\infty} \int_{c}^{R} \frac{1}{x}\,dx
= \underbrace{\lim_{\varepsilon_2\to 0^+} \left(\ln c - \ln \varepsilon_2\right)}_{+\infty}
+ \underbrace{\lim_{R\to\infty} \left(\ln R - \ln c\right)}_{+\infty}.
\label{hauptwert:equation:rechts}
\end{equation}
Beide Grenzwerte divergieren gegen $+\infty$,
das rechte Teilintegral existiert also nicht.
Aus Symmetriegründen divergiert das linke Teilintegral
mit den Grenzwerten $\varepsilon_1 \to 0^+$ und $R' \to -\infty$
genauso gegen $-\infty$.
Insgesamt treten also vier voneinander unabhängige Grenzwerte auf.
Die Aufteilung \eqref{hauptwert:equation:aufteilung} führt also auf den
unbestimmten Ausdruck $\infty - \infty$,
das uneigentliche Integral existiert nicht.

\subsection{Der Cauchy-Hauptwert --- Symmetrische Annäherung}
\input{papers/hauptwert/fig/kehrwert-symmetrisch.tex}
Die zugrunde liegende Idee vom Cauchy-Hauptwert ist, dass
man sich der Polstelle symmetrisch annähert, das heisst mit
einem einzigen, gemeinsamen $\varepsilon$ und symmetrischen Grenzen $\pm R$.
Von den vier unabhängigen Grenzwerten bleiben so nur noch zwei übrig.
Abbildung \ref{hauptwert:figure:kehrwert-symmetrisch} zeigt diese symmetrische
Wahl der Integrationsgrenzen im Vergleich zur unsymmetrischen Aufteilung
aus Abbildung \ref{hauptwert:figure:kehrwert}. Durch symmetrische Annäherung
--- man schreibt dann oftmals $\pv$ (Principal Value) --- erhält man
\begin{equation}
\begin{aligned}
\pv\int_{-\infty}^{\infty} \frac{1}{x}\,dx
&= \lim_{\substack{\varepsilon\to 0^+\\ R\to\infty}}
\biggl(
\int_{-R}^{-\varepsilon} \frac{1}{x}\,dx + \int_{\varepsilon}^{R} \frac{1}{x}\,dx
\biggr) \\
&= \lim_{\substack{\varepsilon\to 0^+\\ R\to\infty}}
\Bigl(
\bigl(\ln\varepsilon - \ln R\bigr) + \bigl(\ln R - \ln\varepsilon\bigr)
\Bigr) \\
&= 0.
\end{aligned}
\label{hauptwert:equation:symmetrisch}
\end{equation}
Der Grenzwert $\varepsilon \to 0^+$, $R \to \infty$ existiert somit und ist $0$,
im Gegensatz zum unabhängigen Grenzwert aus
\eqref{hauptwert:equation:aufteilung}, welcher nicht existiert.

\subsubsection{Was passiert, wenn nicht symmetrisch angenähert wird?}
Eine systematische, aber nicht symmetrische Annäherung bringt auch ein
Resultat, jedoch ist es falsch. Nehmen wir zum Beispiel die verschiedenen
Integrationsgrenzen $\varepsilon$ und $2\varepsilon$. Der Einfachheit halber
wird hier über $[-1, 1]$ integriert, statt über den Bereich $(-\infty, \infty)$.
Aus denselben Symmetriegründen wie in
\eqref{hauptwert:equation:symmetrisch} ist $\pv\int_{-1}^{1} \frac{dx}{x} = 0$.
Mit der unsymmetrischen Annäherung erhalten wir dagegen
\begin{equation}
  \begin{aligned}
    \lim_{\varepsilon \to 0^+}
    \biggl(
    \int_{-1}^{-\varepsilon} \frac{1}{x}\,dx + \int_{2\varepsilon}^{1} \frac{1}{x}\,dx
    \biggr)
    &= \lim_{\varepsilon \to 0^+}
       \bigl((\ln\varepsilon - \ln 1) + (\ln 1 - \ln 2\varepsilon)\bigr) \\
    &= \lim_{\varepsilon \to 0^+} (\ln\varepsilon - \ln 2 - \ln\varepsilon) \\
    &= -\ln 2,
  \end{aligned}
\end{equation}
also nicht den Hauptwert $0$.
Der Grenzwert existiert zwar, er hängt aber vom Verhältnis der beiden
Annäherungen ab und ist damit kein Wert, den man dem Integral zuschreiben kann.

Dies ist die Vorgehensweise des Hauptwerts im Reellen.
Da wir uns in $\mathbb{R}$ auf einer eindimensionalen Zahlengeraden bewegen
und die Integrationswege stets entlang dieser Geraden verlaufen,
bleibt uns nichts anderes übrig, als die Polstelle zu überspringen
und uns ihr symmetrisch anzunähern.
Wie in Abschnitt~\ref{hauptwert:subsection:halbkreiskontur} gezeigt wird,
haben wir diese Einschränkung im Komplexen nicht.
Die komplexen Zahlen bilden nicht nur eine Gerade, sondern eine ganze Ebene.
Dort können wir die reelle Zahlengerade verlassen und einen kleinen Umweg
um die Polstelle machen, nämlich einen Halbkreis
(Abbildung~\ref{hauptwert:figure:halbkreiskontur}).

\subsection{Definition}
Die symmetrische Annäherung lässt sich für beliebige Integranden
formalisieren.
Es gibt zwei Fälle zu unterscheiden. Der erste Fall
(Definition~\ref{hauptwert:definition:pol}) betrifft eine
Divergenz in der Mitte des Integranden in Form einer Singularität.
Der zweite Fall (Definition~\ref{hauptwert:definition:unendlich}) betrifft eine
Divergenz an den unendlichen Integrationsgrenzen $\pm\infty$.

\begin{definition}
\label{hauptwert:definition:pol}
Sei $f$ auf $[a,b]\setminus\{x_0\}$ stetig mit einer Singularität
an der inneren Stelle $x_0 \in (a,b)$.
Der \emph{Cauchy-Hauptwert} des Integrals von $f$ über $[a,b]$ ist
\begin{equation*}
\pv\int_a^b f(x)\,dx
= \lim_{\varepsilon\to 0^+}
\biggl(
\int_a^{x_0-\varepsilon} f(x)\,dx + \int_{x_0+\varepsilon}^{b} f(x)\,dx
\biggr),
\end{equation*}
sofern dieser Grenzwert existiert
\cite{hauptwert:legua-sanchezruiz}.
\end{definition}

\begin{definition}
\label{hauptwert:definition:unendlich}
Sei $f$ auf $\mathbb{R}$ stetig.
Der \emph{Cauchy-Hauptwert} des Integrals von $f$ über $\mathbb{R}$ ist
\begin{equation*}
\pv\int_{-\infty}^{\infty} f(x)\,dx
= \lim_{R\to\infty} \int_{-R}^{R} f(x)\,dx,
\end{equation*}
sofern dieser Grenzwert existiert
\cite{hauptwert:legua-sanchezruiz}.
\end{definition}

Treten wie bei $f(x) = 1/x$ beide Arten der Divergenz gleichzeitig auf,
werden beide Regularisierungen kombiniert:
\begin{equation}
\pv\int_{-\infty}^{\infty} \frac{1}{x}\,dx
= \lim_{\substack{\varepsilon\to 0^+\\ R\to\infty}}
\biggl(
\int_{-R}^{-\varepsilon} \frac{dx}{x} + \int_{\varepsilon}^{R} \frac{dx}{x}
\biggr)
= 0.
\label{hauptwert:equation:kombiniert}
\end{equation}
